% =====================================================================
% 5. Factor models and maximum-likelihood estimation -- question (iv)
% =====================================================================
\section{Factor Models and Estimation}
\label{sec:factormodels}

\subsection{M1: EMPIRICAL DISTRIBUTION}
Under M1, $\Theta$ takes each observed pair $\theta_t=(\theta_{t1},\theta_{t2})'$, $t=1,\dots,n$, with probability $1/n$. Nothing is estimated. Pairs are kept together, so the empirical dependence between the two indices is preserved. The moments of this distribution are the sample mean and the sample covariance with divisor $n$.

\subsection{M2: BIVARIATE GAUSSIAN DISTRIBUTION}
For an i.i.d.\ sample from the bivariate Gaussian distribution with mean $\mu$ and covariance matrix $\Sigma$, the log-likelihood is
\[
\ell(\mu,\Sigma)=-n\log(2\pi)-\frac n2\log|\Sigma|-\frac12\sum_{t=1}^n(\theta_t-\mu)'\Sigma^{-1}(\theta_t-\mu).
\]
Setting $\partial\ell/\partial\mu=\Sigma^{-1}\sum_t(\theta_t-\mu)=0$ gives $\hat\mu=\bar\theta=n^{-1}\sum_t\theta_t$. Substituting $\hat\mu$ and writing $S=n^{-1}\sum_t(\theta_t-\bar\theta)(\theta_t-\bar\theta)'$, the profile log-likelihood is $-\frac n2\{\log|\Sigma|+\mathrm{tr}(\Sigma^{-1}S)\}$ up to a constant. Differentiating with respect to $\Sigma^{-1}$, using $\partial\log|\Sigma^{-1}|/\partial\Sigma^{-1}=\Sigma$ and $\partial\,\mathrm{tr}(\Sigma^{-1}S)/\partial\Sigma^{-1}=S$, gives $\Sigma-S=0$. The maximum-likelihood estimators are thus $\hat\mu=\bar\theta$ and $\hat\Sigma=S$, that is
\[
\hat\sigma_i^2=\frac1n\sum_{t=1}^n(\theta_{ti}-\bar\theta_i)^2,\quad
\hat\rho=\frac{\frac1n\sum_{t=1}^n(\theta_{t1}-\bar\theta_1)(\theta_{t2}-\bar\theta_2)}{\hat\sigma_1\hat\sigma_2}.
\]
M1 and M2 consequently have the same mean vector and covariance matrix, and by (\ref{eq:sk}) the same $s_k$. Asymptotic standard errors are $\hat\sigma_i/\sqrt n$ for $\hat\mu_i$, $\hat\sigma_i/\sqrt{2n}$ for $\hat\sigma_i$ and $(1-\hat\rho^2)/\sqrt n$ for $\hat\rho$.

\subsection{M3: GAUSSIAN MARGINS AND $t$-COPULA}
Let $z_i=(\theta_i-\mu_i)/\sigma_i$ and $u_i=\Phi(z_i)$, and let $\varphi$ be the standard normal density. By Sklar's theorem the joint density of $\Theta$ is the product of the marginal densities and of the copula density evaluated at $(u_1,u_2)$ \citep[Sect.~5.1]{McNeil05}:
\begin{equation}
\label{eq:m3dens}
f_\Theta(\theta_1,\theta_2)=c_{\nu,\rho}(u_1,u_2)\,\prod_{i=1}^2\frac{1}{\sigma_i}\,\varphi(z_i).
\end{equation}
The $t$-copula is the copula of the standard bivariate $t$ distribution with $\nu$ degrees of freedom and correlation parameter $\rho$ \citep[Ch.~5]{McNeil05}, $C_{\nu,\rho}(u_1,u_2)=F_{t,2}\big(t_\nu^{-1}(u_1),t_\nu^{-1}(u_2);\nu,\rho\big)$, where $F_{t,2}$ is the distribution function belonging to the density $f_{t,2}$ and $t_\nu$ is the univariate $t$ distribution function with density $f_{t,1}$. Differentiating with respect to $u_1$ and $u_2$, and using $\partial t_\nu^{-1}(u)/\partial u=1/f_{t,1}(t_\nu^{-1}(u);\nu)$, gives its density
\begin{equation}
\label{eq:tcopdens}
c_{\nu,\rho}(u_1,u_2)=\frac{f_{t,2}(x_1,x_2;\nu,\rho)}{f_{t,1}(x_1;\nu)\,f_{t,1}(x_2;\nu)},\quad x_i=t_\nu^{-1}(u_i)=t_\nu^{-1}\big(\Phi(z_i)\big).
\end{equation}
The log-likelihood of the six parameters $\psi=(\mu_1,\mu_2,\sigma_1,\sigma_2,\rho,\nu)$ is therefore
\begin{equation}
\label{eq:m3lik}
\begin{split}
\ell(\psi)=\sum_{t=1}^n\Big\{&\sum_{i=1}^2\Big[-\log\sigma_i-\tfrac12\log(2\pi)-\tfrac12z_{ti}^2\Big]\\
&+\log f_{t,2}(x_{t1},x_{t2};\nu,\rho)-\sum_{i=1}^2\log f_{t,1}(x_{ti};\nu)\Big\},
\end{split}
\end{equation}
with the standard bivariate and univariate $t$ log-densities
\begin{align*}
\log f_{t,2}(x_1,x_2;\nu,\rho)&=\log\Gamma\big(\tfrac{\nu+2}2\big)-\log\Gamma\big(\tfrac\nu2\big)-\log(\pi\nu)-\tfrac12\log(1-\rho^2)\\
&\quad-\tfrac{\nu+2}{2}\log\Big(1+\frac{x_1^2-2\rho x_1x_2+x_2^2}{\nu(1-\rho^2)}\Big),\\
\log f_{t,1}(x;\nu)&=\log\Gamma\big(\tfrac{\nu+1}2\big)-\log\Gamma\big(\tfrac\nu2\big)-\tfrac12\log(\pi\nu)-\tfrac{\nu+1}2\log\Big(1+\frac{x^2}\nu\Big).
\end{align*}
Unlike in M2, the marginal parameters also enter the copula term through $x_{ti}$, so the first-order conditions have no closed-form solution and (\ref{eq:m3lik}) is maximised numerically. We proceed in two steps. First, the margins are estimated by their own Gaussian likelihood and $(\rho,\nu)$ by maximising the copula term alone; this is the inference-functions-for-margins estimator \citep[Sect.~5.5]{McNeil05}. Second, starting from this point, all six parameters are optimised jointly on (\ref{eq:m3lik}) with a quasi-Newton method (L-BFGS-B) under the constraints $|\rho|<1$ and $2<\nu\le 200$, with $\log\nu$ as the optimisation variable. Only the second step delivers the maximum-likelihood estimator of the full model, which is efficient and recognises that the margins also shape the copula term; the first provides its starting value and a robustness check. Standard errors are obtained from the inverse of the observed information matrix, computed by finite differences. For $\nu\to\infty$ the $t$-copula converges to the Gaussian copula and M3 to M2, so that M2 is nested in M3 on the boundary of the parameter space.

Two quantities implied by M3 are needed later. The coefficient of lower (and, by symmetry, upper) tail dependence of the $t$-copula is \citep[Example~5.33]{McNeil05}
\begin{equation}
\label{eq:taildep}
\lambda_L=\lim_{q\downarrow0}\P[U_2\le q\mid U_1\le q]=2\,t_{\nu+1}\!\left(-\sqrt{\frac{(\nu+1)(1-\rho)}{1+\rho}}\right),
\end{equation}
whereas $\lambda_L=0$ for the Gaussian copula with $\rho<1$. The Pearson correlation of $\Theta$ under M3 is not equal to $\rho$ and has no closed form; it is computed by numerical integration of (\ref{eq:m3dens}) and used in (\ref{eq:sk}).

\subsection{ESTIMATES}
\begin{table}[tb]
\caption{Maximum-likelihood estimates of M2 and M3 (standard errors in parentheses).}
\label{tab:mle}
\centering
\qrmtable{tab_mle}
\end{table}

Table~\ref{tab:mle} reports the estimates. The marginal parameters of the two models almost coincide. The $t$-copula has $\hat\nu=\val{q4.M3.params.5|f1}$ degrees of freedom with a standard error of \val{q4.M3.se.5|f1}, and implies a tail-dependence coefficient of \val{q4.M3.tail_dep|f3}. The likelihood-ratio statistic for the Gaussian copula against the $t$-copula is $2(\ell_{M3}-\ell_{M2})=\val{q4.LR|f1}$. Because the null hypothesis $\nu=\infty$ lies on the boundary of the parameter space, the statistic is asymptotically distributed as an equal mixture of a point mass at zero and a $\chi^2_1$ \citep{SelfLiang87}; the $p$-value is below $10^{-9}$ and the Gaussian copula is rejected. The first-step estimates, $\hat\rho=\val{q4.M3.ifm.4|f3}$ and $\hat\nu=\val{q4.M3.ifm.5|f1}$, are very close to the joint ones.

Rejecting M2 against M3 does not make M3 adequate. M3 keeps Gaussian margins, and Table~\ref{tab:desc} and Figure~\ref{fig:qq} show that the margins are where the data depart most from normality. Figure~\ref{fig:scatter} makes the point visually: samples from M2 and M3 look alike and neither reproduces the crash weeks in the lower-left corner of the data.

\begin{figure}[tb]
\centering
\qrmfigure{fig_factors_scatter}
\caption{Observed weekly log-returns (left) and samples of the same size from the fitted M2 and M3.}
\label{fig:scatter}
\end{figure}
